\section{BOLSIG+ Governing Equations}
The electron Boltzmann equation may be written as follows:
%
\begin{equation}
\pp{f}{t}
+
\mbf{v} \cdot \nabla f
-
\frac{e}{m} \mbf{E} \cdot \nabla_{\mbf{v}} f
=
C[f],
\label{eqn:electron-boltzmann-full}
\end{equation}
%
where $f(\mbf{x}, \mbf{v}, t)$ denotes the electron distribution
function, $\mbf{v}$ is the velocity vector, $\nabla f$ is the spatial
gradient of $f$ (i.e., wrt $\mbf{x}$), $\nabla_{\mbf{v}} f$ is the
velocity gradient of $f$ (i.e., wrt $\mbf{v}$), $e$ is the absolute
value of the charge of an electron, $m$ is the mass of an electron,
$\mbf{E}$ is the electric field, and $C[f]$ denotes the collision
term.

The development here is based upon the description of Hagelaar and
Pitchford~\cite{} (referred to as HP for short) and the BOLSIG+ User's
Manual~\cite{}.

To begin, it is assumed that the electric field and collision
probabilities are spatially uniform.  In this situation, HP asserts
that the distribution function $f$ is 1) ``symmetric in velocity space
around the electric field direction'' and 2) ``may vary only along the
field direction'' in physical space.  The meaning of 1) is somewhat
unclear, but in the subsequent analysis, a spherical coordinate system
is used and no possible dependence on the azimuthal angle is included.
Thus, the combined import of these assumptions is to reduce the
dependence of $f$ from 6D plus time to 3D plus time.  Specifically,
letting $z$ denote the direction of the electric field, $v$ denote the
velocity magnitude, and $\theta$ denote the polar angle between the
electric field direction and the velocity direction, we have $f = f(z,
v, \theta, t)$.  In this situation,~\ref{eqn:electron-boltzmann-full}
simplifies to
%
\begin{equation}
\pp{f}{t}
+
v \cos \theta \pp{f}{z}
-
\frac{e}{m} E \left( \cos \theta \pp{f}{v} - \frac{\sin \theta}{v} \pp{f}{\theta} \right)
=
C[f].
\label{eqn:electron-boltzmann-reduced}
\end{equation}
%

In preparation for subsequent development, let $\xi = \cos \theta$.
Then,~\ref{eqn:electron-boltzmann-reduced} becomes
%
\begin{equation}
\pp{f}{t}
+
v \xi \pp{f}{z}
-
\frac{e}{m} E \left( \xi \pp{f}{v} + \frac{(1-\xi^2)}{v} \pp{f}{\xi} \right)
=
C[f].
\label{eqn:electron-boltzmann-xi}
\end{equation}
%
To continue, the following two-term expansion is introduced:
%
\begin{equation}
f(v, \xi, z, t) = f_0(v, z, t) + \xi f_1(v, z, t).
\label{eqn:two-term}
\end{equation}
%
To derive governing equations for $f_0$ and $f_1$, we apply the
Galerkin method in the $\xi$ coordinate.  Specifically,
substitute~\eqref{eqn:two-term} into
~\eqref{eqn:electron-boltzmann-xi} and integrate against $1$ and
$\xi$.  Rewriting the result in terms of a rescaled energy variable
$\epsilon \equiv (v^2 / \gamma^2)$, where $\gamma^2 = 2e / m$, gives
%
\begin{gather}
\pp{f_0}{t}
+
\frac{\gamma}{3} \sqrt{\epsilon} \pp{f_1}{z}
-
\frac{\gamma}{3} \frac{E}{\sqrt{\epsilon}} \pp{}{\epsilon} \left( \epsilon f_1 \right)
=
C_0
\label{eqn:f0}\\
%
\pp{f_1}{t}
+
\gamma \sqrt{\epsilon} \pp{f_0}{z}
-
\gamma E \sqrt{\epsilon} \pp{f_0}{\epsilon}
=
C_1
\label{eqn:f1}
\end{gather}
%
where
%
\begin{gather*}
C_0 = \frac{1}{2} \int_{-1}^{1} C[f_0 + \xi f_1] \, d\xi\\
C_1 = \frac{3}{2} \int_{-1}^{1} \xi C[f_0 + \xi f_1] \, d\xi.
\end{gather*}
%
Note that~\eqref{eqn:f0} and~\eqref{eqn:f1} are equivalent to HP (5)
and (6), respectively.  Although, to match the forms in HP more
precisely, we must examine the collision terms in more detail.
\todo[inline]{Collisions}
